\section{İlgili Çalışmalar}
\label{sec:related}

Bu çalışma, beş ayrı fakat birbiriyle yakından ilişkili literatür
çizgisiyle bağlantılıdır.

\subsection{Kompakt Yinelemeli Hücreler}

Standart LSTM~\cite{hochreiter1997lstm} ve GRU~\cite{cho2014gru}
hücreleri uzun süreli zamansal örüntüleri etkili biçimde yakalar;
ancak kilobayt düzeyinde belleğe sahip MCU'lar için ağır kalır.
Örneğin orta büyüklükte bir LSTM ($H{=}64$, $d{=}3$) FP32 duyarlıkta
daha şimdiden \SI{50}{\kilo\byte} üzerinde ağırlık gerektirir. Bu
değer, hedef aygıtımızın \SI{16}{\kilo\byte} Flash sınırını aşar.

\fastgrnn{}~\cite{kusupati2018fastgrnn}, bu boşluğu düşük ranklı
ağırlık matrisleriyle desteklenen iki skaler kapılı ($\zeta, \nu$)
bir hücreyle kapatmayı amaçlar. Dizi görevlerinde LSTM ile
karşılaştırılabilir doğruluğa, bir ila iki büyüklük mertebesi daha
az parametreyle ulaşabilir. Yinelemeli olmayan tamamlayıcı EdgeML
yaklaşımları arasında sığ doğrusal olmayan bir ağaç olan
Bonsai~\cite{kumar2017bonsai} ve prototip tabanlı bir sınıflandırıcı
olan ProtoNN~\cite{gupta2017protonn} yer alır. Bu üç yöntem
Microsoft Research EdgeML araç takımını oluşturur~\cite{edgeml}.
HAR için kritik olan zamansal girdiyi, bu üçlü içinde doğrudan
işleyen yöntem \fastgrnn{}'dur.

\subsection{Nicelendirme}

Eğitim-sonrası nicelendirme (PTQ), sinir ağı ağırlıklarını ve
aktivasyonlarını FP32 yerine sabit noktalı ya da düşük bitli tamsayı
biçimlerinde temsil eder. Jacob \emph{vd.}~\cite{jacob2018integer},
tensör-başına ölçek kullanan ve yalnızca tamsayı aritmetiğine dayanan
çıkarımı tanıtmıştır. Han \emph{vd.}~\cite{han2016deepcompression}
ise budama, nicelendirme ve Huffman kodlamasını birleştirerek derin
ağları büyük ölçüde sıkıştırmıştır. Nicelendirme farkında eğitim
(QAT)~\cite{hubara2017qnn}, PTQ sonrasında görülebilen doğruluk
kaybını azaltır; ancak bunun bedeli yeniden eğitimdir. Bu çalışmada
hem ağırlıklar hem de, açık aktivasyon kalibrasyonuyla, ara tensörler
için tensör-başına Q15 PTQ kullanıyoruz. Kalibrasyon adımı pratikte
kritiktir: Bölüm~\ref{sec:results:quant}'da gösterildiği gibi,
kalibrasyon yapılmadığında model çöker; kalibrasyon yapıldığında ise
QAT'ye gerek kalmadan doğruluk korunur.

\subsection{Seyreklik ve Budama}

\fastgrnn{}'da kullanılan seyrekleştirme yöntemi iteratif sert
eşiklemedir (IHT)~\cite{blumensath2009iht}. Her eğitim adımında,
her ağırlık tensöründeki mutlak değerce en büyük $k$ girdi korunur;
diğer girdiler sıfırlanır ve ağ bu sabit maske ile eğitilmeye devam
eder. Han \emph{vd.}~\cite{han2015pruning}, daha büyük ağlarda
agresif budamanın doğruluğu koruyabildiğini göstermiştir. Frankle ve
Carbin'in piyango bileti hipotezi~\cite{frankle2019lt} de bu tür
seyrek alt ağların neden eğitilebilir kaldığına kuramsal bir zemin
sağlar.

\subsection{Daha Güçlü Hedeflerde tinyML}

tinyML literatürü~\cite{banbury2021mlperf} çoğunlukla ARM Cortex-M
sınıfı hedeflere odaklanır. Bu hedeflerde yüzlerce kilobayt SRAM,
donanım FPU'ları ve tek çevrimli çarpıcılar bulunur. MCUNet~\cite{lin2020mcunet},
Cortex-M verimini artırmak için sinir mimarisini ve yorumlayıcıyı
birlikte tasarlar. CMix-NN~\cite{capotondi2020cmixnn}, bellek
kısıtlı uç aygıtlar için karma duyarlıklı çekirdekler sunar.
ARM'ın CMSIS-NN kütüphanesi~\cite{lai2018cmsisnn}, yaygın katmanları
Cortex-M SIMD yönergeleriyle hızlandırır. TensorFlow Lite
Micro~\cite{david2021tflitemicro} ise taşınabilir bir çalışma
zamanı katmanı sağlar.

Bu çalışmaların hiçbiri, bu makalede ele alınan \msp{} sınıfı aygıtı
hedeflemez. Bu aygıt bir ila iki büyüklük mertebesi daha küçüktür ve
donanım çarpıcısına sahip değildir. Bildiğimiz kadarıyla bu çalışma,
donanım çarpıcısı bulunmayan bir MCU üzerinde kapılı yinelemeli bir
hücrenin herkese açık ilk uçtan uca yeniden üretimidir.

\subsection{İnsan Aktivite Tanıma}

UCI HAR~\cite{anguita2013public} ve geçiş etiketlerini de içeren
uzantısı \hapt{}~\cite{reyes2015transition}, ivmeölçer tabanlı HAR
için yaygın kullanılan kamusal kıyaslama kümeleridir.
DeepConvLSTM~\cite{ordonez2016deepconvlstm}, yığılmış CNN ve LSTM
katmanlarıyla güçlü bir derin öğrenme taban çizgisi oluşturmuştur.
Wang \emph{vd.}~\cite{wang2019harsurvey} ile Demrozi
\emph{vd.}~\cite{demrozi2020harsurvey} tarafından hazırlanan
taramalar, yöntemsel manzarayı geniş biçimde özetler. Her iki
tarama da dinamik \textsc{downstairs} sınıfının kalıcı bir zorluk
olduğunu vurgular. Bölüm~\ref{sec:results:perclass}'ta bu bulgunun
bizim deneylerimizde de sürdüğünü gösteriyoruz.
